% BHU PYQ -- Manually transcribed & error-corrected
% Paper: BCH-223 : Specialised Accounts-II
% Exam: B.Com. (Hons.) Semester IV Examination, 2023-24

\documentclass[11pt,a4paper]{article}
\usepackage[a4paper,top=2.5cm,bottom=2.5cm,left=2.8cm,right=2.8cm,headheight=14pt]{geometry}
\usepackage{amsmath,amssymb}
\usepackage[T1]{fontenc}
\usepackage[utf8]{inputenc}
\usepackage{lmodern,microtype,fancyhdr,enumitem,hyperref,lastpage,parskip}

\hypersetup{
    colorlinks=true,
    urlcolor=black,
    linkcolor=black,
    pdftitle={BCH-223: Specialised Accounts-II -- BHU B.Com. (Hons.) Sem IV 2023-24}
}

\pagestyle{fancy}
\fancyhf{}
\renewcommand{\headrulewidth}{0.4pt}
\renewcommand{\footrulewidth}{0.4pt}
\fancyhead[L]{\small   \textit{Banaras Hindu University}}
\fancyhead[R]{\small   \textit{BCH-223: Specialised Accounts-II}}
\fancyfoot[C]{\small Page  \thepage\ of \pageref{LastPage}}


\newlist{parts}{enumerate}{2}
\setlist[parts,1]{label=(\roman*),leftmargin=*,itemsep=5pt,topsep=3pt}
\setlist[parts,2]{label=(\alph*),leftmargin=*,itemsep=3pt,topsep=2pt}

\newcommand{\pts}[1]{\hfill   \textit{[#1]}}


\begin{document}


\begin{center}
  {\large\bfseries BANARAS HINDU UNIVERSITY}\\[2pt]
  {\normalsize B.Com. (Hons.) --- Semester IV Examination, 2023-24}\\[6pt]
  \rule{0.8\linewidth}{0.5pt}\\[6pt]
  {\large\bfseries Paper No. BCH-223: Specialised Accounts-II}\\[4pt]
  {\normalsize Time: 3 Hours\hfill Maximum Marks: 70}
\end{center}

\rule{\linewidth}{0.4pt}

\smallskip



\noindent   \textit{Instructions: Answer all questions. All questions carry equal marks (14 marks each).}

\bigskip




\noindent   \textbf{Question 1.}
Explain the various schedules which are prepared for Balance Sheet of Bank and give its specimen. What is slip system? Explain and illustrate.\pts{14}

\centerline{   \textbf{OR}}

The following balances appeared in the books of   \textit{Bank of Prayagraj} as on 31st March, 2024 (amounts in thousands, i.e., Rs. 000):


\begin{center}
\small

\begin{tabular}{|l|r||l|r|}
\hline
   \textbf{Particulars} &   \textbf{Amount (Rs. )} &   \textbf{Particulars} &   \textbf{Amount (Rs. )} \\
\hline
Share Capital & 9,000 & Gold & 1,000 \\
Reserve Fund & 1,375 & Silver & 757 \\
Fixed Deposits & 34,667 & Investments in Securities at Cost & 64,070 \\
Current Accounts & 64,758 & \quad (Market Price: Rs. 65,250) & \\
Savings Bank Account & 20,002 & Interest Accrued on Investments & 452 \\
Cash in Hand & 8,178 & Cash Credits, Overdrafts \& Demand Loans & 33,719 \\
Balance with RBI & 503 & Term Loans & 6,505 \\
Balances with Banks & 12,995 & Bills Purchased and Discounted & 13,714 \\
Debts due to Bank and Agents & 7,741 & Furniture & 300 \\
Branch Adjustment (Cr.) & 4,499 & Vehicles & 310 \\
Rebate on Bills Discounted & 100 & Money at Call and Short Notice & 525 \\
Acceptances for Customers & 5,747 & Profit and Loss Account (Cr.) & 976 \\
Bills for Collection & 3,353 & \quad (incl. current year's profit of Rs. 900) & \\
\hline
\end{tabular}
\end{center}

Prepare the Balance Sheet as on 31st March, 2024 after taking the following information into consideration:

\begin{parts}
  \item The authorized capital of the bank is Rs. 2.5 crores of rupees in 2,50,000 shares of Rs. 100 each, of which 2,00,000 shares are issued at Rs. 45 per share being called and paid up.
  \item Claims against the bank not acknowledged as debts amount to Rs. 1,73,000.
\end{parts}\pts{14}

\medskip\rule{\linewidth}{0.2pt}

\pagebreak




\noindent   \textbf{Question 2.}
State briefly the principles of Insurance and also prepare Revenue Account and Balance Sheet of a General Insurance Company with imaginary figures. What is Double Insurance? How is it different from Re-Insurance?\pts{14}

\centerline{   \textbf{OR}}

From the following particulars you are required to prepare Fire Insurance Revenue Account and Profit and Loss Account for the year ending on 31st March, 2024 (amounts in thousands, i.e., Rs. 000):


\begin{center}
\small

\begin{tabular}{|l|r|}
\hline
   \textbf{Particulars} &   \textbf{Amount (Rs. )} \\
\hline
Premium received & 6,00,000 \\
Re-insurance premium & 60,000 \\
Claims paid & 2,40,000 \\
Claims outstanding on 1st April, 2023 & 20,000 \\
Claims intimated but not accepted on 31st March, 2024 & 5,000 \\
Claims intimated and accepted but not paid on 31st March, 2024 & 30,000 \\
Commission & 1,00,000 \\
Commission on reinsurance ceded & 4,000 \\
Commission on reinsurance accepted & 2,000 \\
Expenses of Management & 1,01,000 \\
Provision for unexpired risks on 1st April, 2023 & 2,00,000 \\
Additional provision for unexpired risks on 1st April, 2023 & 10,000 \\
Re-insurance recoveries of claims & 4,000 \\
Legal expenses regarding claims & 2,500 \\
Loss on sale of Motor Car & 1,750 \\
Bad Debts & 1,250 \\
Interest and Dividend (Gross) & 45,000 \\
Income Tax Paid & 9,000 \\
Rent of staff quarters deducted from salaries & 700 \\
Depreciation of Furniture & 2,800 \\
\hline
\end{tabular}
\end{center}

You are required to provide for additional provision for unexpected (unexpired) risks at 1\% of the net premium in addition to the opening balance.\pts{14}

\medskip\rule{\linewidth}{0.2pt}

\pagebreak




\noindent   \textbf{Question 3.}
From the following particulars, prepare Statement of Affairs and Deficiency Account of Mr. A of Mumbai:

\begin{itemize}
  \item Bank overdraft Rs. 1,00,000 secured against the mortgage of land;
  \item Trade creditors Rs. 60,000;
  \item Bills discounted are Rs. 10,000 out of which Rs. 4,000 are doubtful;
  \item His private liabilities are Rs. 10,000;
  \item Salary due to two clerks Rs. 1,200;
  \item Wages due to a servant Rs. 500;
  \item Landlord's two months' rent due Rs. 400;
  \item Manager's salary due Rs. 900;
  \item Mr. A is a trustee of his minor nephew valued at Rs. 5,000;
  \item His household assets are valued at Rs. 4,000;
  \item His wife's ornaments are of Rs. 10,000, out of which ornaments of Rs. 4,000 have been purchased from the drawings of her husband and the remaining from the amount received from her father.
\end{itemize}
His business assets are: Cash Rs. 500, Furniture Rs. 3,000 (realisable Rs. 2,000), Land Rs. 60,000 (realisable Rs. 40,000), Machinery Rs. 70,000 (realisable Rs. 50,000), Stock Rs. 5,000 (realisable Rs. 2,000), Book Debts Rs. 20,000 (realisable Rs. 12,000), Business capital Rs. 45,000, Losses Rs. 49,500.\pts{14}

\centerline{   \textbf{OR}}

Mr. X of Delhi was declared insolvent on 31st December, 2023 when his positions were as under:


\begin{center}
\small

\begin{tabular}{|l|r|r|}
\hline
   \textbf{Particulars} &   \textbf{Debit (Rs. )} &   \textbf{Credit (Rs. )} \\
\hline
Cash & 1,000 & -- \\
Furniture (estimated to produce Rs. 2,000) & 5,000 & -- \\
Stock (estimated to produce Rs. 1,000) & 3,000 & -- \\
Shares (estimated to produce Rs. 1,000) & 3,000 & -- \\
Machinery (estimated to produce Rs. 13,000) & 20,000 & -- \\
Sundry Debtors: & & \\
\quad -- Good & 2,000 & -- \\
\quad -- Doubtful (estimated to produce 25\%) & 4,000 & -- \\
\quad -- Bad & 2,000 & -- \\
Loss (2021 and 2022) & 3,000 & -- \\
Profit (2023) & -- & 1,000 \\
Capital & -- & 5,000 \\
Sundry Creditors: & & \\
\quad -- Trade Creditors & -- & 12,500 \\
\quad -- Bills Payable & -- & 3,000 \\
\quad -- Bank Overdraft & -- & 12,000 \\
\quad -- Wife Smt. Z & -- & 4,000 \\
\quad -- Father Mr. Y & -- & 1,200 \\
\quad -- Preferential Creditors & -- & 300 \\
Profit through betting & -- & 4,000 \\
\hline
   \textbf{Total} &   \textbf{43,000} &   \textbf{43,000} \\
\hline
\end{tabular}
\end{center}

The insolvent has mortgaged machinery with the bank. Smt. Z lent Rs. 3,000 out of her `Stridhan'. Shares have been given as security to Mr. Y. The business was started on 1st January, 2021. He had transferred some property to his wife in the following manner:

\begin{parts}
  \item Property worth Rs. 1,000 transferred on 1st July, 2023 without consideration.
  \item Property worth Rs. 2,000 transferred on 1st July, 2022 for a sufficient consideration.
  \item Property worth Rs. 2,500 transferred on 1st July, 2021 without consideration.
\end{parts}
Prepare Mr. X's Statement of Affairs and Deficiency Account.\pts{14}

\medskip\rule{\linewidth}{0.2pt}

\pagebreak




\noindent   \textbf{Question 4.}
Discuss the concept and features of Double Account System. What is the valuable purpose which it serves? How are final accounts prepared in it? Give example.\pts{14}

\centerline{   \textbf{OR}}

Following is the Trial Balance of ABC Company Ltd. as on 31st March, 2024:


\begin{center}
\small

\begin{tabular}{|l|r|r|}
\hline
   \textbf{Particulars} &   \textbf{Debit (Rs. )} &   \textbf{Credit (Rs. )} \\
\hline
Equity Share Capital & -- & 60,00,000 \\
Preference Share Capital & -- & 15,00,000 \\
Debentures & -- & 40,00,000 \\
Depreciation Fund & -- & 12,00,000 \\
Sundry Creditors & -- & 2,00,000 \\
Transfer Fees & -- & 10,000 \\
Sale of Gas & -- & 42,50,000 \\
Reserve Fund & -- & 10,00,000 \\
Share Premium Reserve & -- & 3,00,000 \\
Sale of Coke & -- & 18,50,000 \\
Balance of Net Revenue A/c (01-04-2023) & -- & 4,50,000 \\
Mains and Meters & 11,50,000 & -- \\
Parliamentary Expenses & 1,02,000 & -- \\
Repair of Office Building & 75,000 & -- \\
Opening Stock of Material & 7,000 & -- \\
Opening Stock of Coal & 4,25,000 & -- \\
Interest on Debentures & 2,30,000 & -- \\
Preliminary Expenses & 4,00,000 & -- \\
Director's and Auditor's Fees & 1,60,000 & -- \\
Land and Buildings & 62,50,000 & -- \\
Wagons and Lorries & 53,000 & -- \\
Purchase of Coal & 22,50,000 & -- \\
Purchase of Materials & 28,000 & -- \\
Legal Charges & 35,000 & -- \\
Dividend Paid on Preference Shares & 14,00,000 & -- \\
Interim Dividend Paid & 6,00,000 & -- \\
Repairs and Renewals of Plants and Stores & 4,50,000 & -- \\
Printing and Stationery & 25,000 & -- \\
Plants and Stores & 47,00,000 & -- \\
Investment of Depreciation Fund & 12,00,000 & -- \\
Sundry Debtors & 5,50,000 & -- \\
Renewal and Maintenance of Mains and Meters & 2,10,000 & -- \\
Compensation to Employees & 50,000 & -- \\
Cash in Hand and Cash at Bank & 3,00,000 & -- \\
Sundry Expenses & 1,10,000 & -- \\
\hline
   \textbf{Total} &   \textbf{2,07,60,000} &   \textbf{2,07,60,000} \\
\hline
\end{tabular}
\end{center}




\noindent   \textbf{Other Information:}

\begin{parts}
  \item Closing stock of Coal and Material as on 31st March, 2024 was Rs. 4,50,000 and Rs. 6,000 respectively.
  \item Reserve fund is to be increased to Rs. 12,50,000.
  \item Depreciation is to be provided as follows: Land and Building Rs. 1,50,000; Mains and Meters Rs. 50,000 and Wagons and Lorries Rs. 1,00,000.
  \item During the year the company expended Rs. 3,00,000 on Land, Rs. 2,00,000 on Mains and Rs. 5,00,000 on Plants.
\end{parts}
On the basis of above information prepare Final Accounts under the Double Account System.\pts{14}

\medskip\rule{\linewidth}{0.2pt}




\noindent   \textbf{Question 5.}
State the meaning, characteristics and the fundamental principles of government accounting. How does it differ from commercial accounting? Explain. What is PAC?\pts{14}

\centerline{   \textbf{OR}}

Write short notes on any two of the following:

\begin{parts}
  \item NPA (Non-Performing Asset)
  \item Profit and Loss Account of Bank
  \item C\&AG (Comptroller and Auditor General)
  \item Deficiency Account
\end{parts}\pts{14}

\vfill
\rule{\linewidth}{0.4pt}

\begin{center}\small
\href{https://bhu-exam.vercel.app/}{Click here to visit our website}\\[4pt]
\href{https://me-aryan.vercel.app/}{Click here to visit our contributor}\\[4pt]
\href{https://quashan-me.vercel.app/}{Click here to visit our contributor}
\end{center}

\end{document}
